\documentclass[12pt]{article}
\usepackage[utf8]{inputenc}
\usepackage[T1]{fontenc}
\usepackage{hyperref}
\title{Yaben Language Resources for K-12 Curriculum Development}
\author{Compiled by Claude Sonnet 5.5 (claude-sonnet-5-5) for Andr\'as Kornai}
\date{}
\begin{document}
\maketitle

\section{Introduction}
Yaben (also Yamben; ISO 639-3 ybm) is a tonal Papuan language of Madang Province, Papua New Guinea, roughly 100 km northwest of Madang town according to Joshua Project. Speaker counts conflict: Wikipedia cites 1{,}090 (2000 census), Joshua Project gives 2{,}400, and Global Recordings Network gives 700. Joshua Project reports that the language is used in everyday conversation across generations, that literacy is low, and that no Scripture translation exists. Educational material is very scarce in my searches.

\section{Primary Educational Resources}
No primers, textbooks or school materials in Yaben turned up in my searches.

\section{Linguistic and Reference Materials}
\begin{itemize}
\item Wikipedia, Yaben language: \url{https://en.wikipedia.org/wiki/Yaben_language}
\end{itemize}
This is an overview page. No dictionary or grammar turned up in my searches.

\section{Religious and Community Resources}
\begin{itemize}
\item Joshua Project, Yaben in Papua New Guinea: \url{https://joshuaproject.net/people_groups/15931}
\end{itemize}
Joshua Project reports that no Scripture exists in the language.

\section{Digital and Audio Resources}
\begin{itemize}
\item Global Recordings Network, Yaben (ybm): \url{https://globalrecordings.net/en/language/ybm}
\end{itemize}
This page was listed in search results, but I did not read its content, so the contents and type of the recordings are unconfirmed.

\section{Curriculum Development Context}
I found no evidence that Yaben is taught in schools or has a standard orthography. A curriculum would likely have to start with documentation and community-led orthography work, and teaching would take place alongside Tok Pisin and English.
\end{document}